\section{Formalization Map}
\label{sec:formalization-map}

\begin{layman}
This final section is the bridge between two parallel tracks. On
one track, the prose of sections 0--8 makes a complete
mathematical argument: a paper-style proof that the Jacobian
construction works. On the other, the Lean code of the
\code{Jacobian/} repository turns each step of that argument into
machine-verified theorems. The bridge below is a simple table
mapping each declaration in \code{Challenge.lean} --- the
public-facing API --- to its informal source: a section, a theorem,
sometimes a chain of lemmas culminating in a single informal
claim.

The same table also lays out where the project still has work to
do. Some informal theorems have been fully formalised; their
table rows are decorated with green check marks (\(\leanok\)) in
the dependency graph. Others are still open, pinned by literal
\code{sorry} in the Lean code; those rows are decorated with
yellow ``not ready'' flags. The largest remaining gaps are the
classical inputs that Mathlib does not yet contain --- the
Riemann--Roch theorem, Stokes on manifolds with boundary, Abel's
theorem, and the surface classification underlying the polygonal
model. Each is a multi-month formalisation project in its own
right, and several are actively in progress in the wider Mathlib
community.

The dependency graph rendered alongside the web blueprint is
generated automatically from these annotations. A reader curious
about a particular theorem can click through its dependencies ---
upstream to the lemmas it uses, downstream to the theorems it
supports --- and trace the entire fabric of the proof until they
reach either an axiom of mathematics or a yet-unfilled box. That
graph is the project's long-term planning instrument: when a
node turns green, every downstream node it was blocking moves one
step closer to formal closure.
\end{layman}

This section maps the informal proof to the current Lean project structure.

\begin{center}
\begin{tabular}{p{0.34\linewidth}p{0.56\linewidth}}
\textbf{Challenge declaration} & \textbf{Informal source} \\
\hline
\code{genus} &
Definition~\ref{def:analytic-genus}; finite-dimensionality from
Input~\ref{input:finite-dimensionality}. \\
\code{genus_eq_zero_iff_homeo} &
Theorem~\ref{thm:genus-zero-classification}; uses Riemann-Roch and
degree-one map classification. \\
\code{Jacobian} and basic instances &
Definition~\ref{def:analytic-jacobian}; period lattice theorem plus complex
torus package. \\
\code{Jacobian.ofCurve} &
Definition~\ref{def:abel-jacobi}; path independence modulo periods. \\
\code{ofCurve_contMDiff} &
Lemma~\ref{lem:aj-holomorphic}. \\
\code{ofCurve_self} &
Lemma~\ref{lem:aj-self}. \\
\code{ofCurve_inj} &
Theorem~\ref{thm:abel-jacobi-injective}; Abel's theorem. \\
\code{pushforward}, \code{pullback} &
Lemma~\ref{lem:push-pull-descend}; trace and pullback preserving periods. \\
\code{pushforward_contMDiff} and \code{pullback_contMDiff} &
Theorem~\ref{thm:push-pull-functoriality}; descended complex-linear maps are
holomorphic. \\
\code{pushforward_id_apply} and \code{pullback_id_apply} &
Theorem~\ref{thm:push-pull-functoriality}; identity laws for pullback and
trace. \\
\code{pushforward_comp_apply} and \code{pullback_comp_apply} &
Theorem~\ref{thm:push-pull-functoriality}; composition laws for pullback and
trace. \\
\code{ContMDiff.degree} &
Definition~\ref{def:degree-trace}. \\
\code{pushforward_pullback} &
Theorem~\ref{thm:pushforward-pullback}; trace-pullback identity. \\
\end{tabular}
\end{center}

The current repo already has production modules for many algebraic and
topological parts of this map. The largest classical gaps are:

\begin{itemize}
\item Riemann-Roch and divisor theory for compact Riemann surfaces.
\item Abel's theorem in point-separation form.
\item Riemann bilinear relations and full-rank period lattice theorem.
\item Branched-cover degree theory and trace of holomorphic differentials.
\item Analytic genus versus topological genus comparison.
\end{itemize}
